% Review
\section{Návrh uživatelského rozhraní}

% Jelikož se aplikace mění ze zaměření na výuku jazyků na výuku softwarového inženýrství a programování, je třeba tomu přizpůsobit i vzhled uživatelského rozhraní. Tomu se budu věnovat v této kapitole.

% Review
\subsection{Výběr barev}

% Jelikož stávající aplikace již vychází ze standardu Material Design 2 a využívá knihovnu Material UI, budu volit zbývající barvy, včetně sekundární barvy, v souladu s tímto standardem.

% Material Design 2 specifikuje tyto základní barvy, které jsou pak konzistentně používány napříč celou aplikací.

% Review
\subsubsection{Primární barva}

% Říct na co se používá primární barva dle https://m2.material.io/design/color/the-color-system.html#color-theme-creation

% Primární barvu budu vybírat podle toho, jakým způsobem chce aplikace působit na uživatele. Vzdělávací aplikace mají za úkol předat uživatelům znalosti a informace.
% Kličový aspekt vyučování - ten co vyučuje je v pozici experta, který by měl garantovat, že jsou vyučované informace pravdivé a kvalitní. Z toho důvodu je naprosto klíčová důvěra uživatele k učiteli. Učitel musí působit profesionálně a důvěryhodně.
% Z tohoto důvodu jsem se rozhodl jako primární barvu zvolit modrou, jelikož modrá barva je považována za barvu důvěryhodnosti, inteligence a přesnosti. (zde přidat odkaz na prezentaci https://docs.google.com/presentation/d/1l_T7H8KyT43RTaGzcsLXo4t1wO8s30aD6yLaMZZ1InQ/edit?slide=id.ga6a0d160f4_0_48#slide=id.ga6a0d160f4_0_48)

% Proto jsem se rozhodl zvolit jako primární barvu modrou #4fb5ea. V rámci Material Design 2 standardu má  

% Review
\subsubsection{Sekundární barva}

% Zde podle toho, jestli mi z designu ve Figmě vyplynulo, že chci sekundární barvu, říct, že buď jsem se rozhodl sekundární barvu nemít kvůli jednoduchosti uživatelského rozhraní (což je perfektně v souladu s Material Design 2 standardem) nebo že jsem se rozhodl sekundární barvu vybrat -> pak říct, proč jsem si vybral komplementární/triadickou/apod. barvu. Když už bych si nějakou vybíral, tak bych si asi vybral pastelovou oranžovou barvu. (viz https://m2.material.io/design/color/the-color-system.html#color-theme-creation)

% Material Design 2 specifikuje kromě primární barvy ještě sekundární barvu. Tato barva je používána zjeména na ... doplnit podle Material UI guidelines a hodit link sem (https://m2.material.io/design/color/the-color-system.html#color-theme-creation)



% Review
\subsubsection{Barvy pozadí}

% Pro normální pozadí použít čistě bílou #ffffff
% Pro dark theme jsem se rozhodl použít barvu tmavě modrou

% Review
\subsubsection{Další barvy}

% Write - zde sepsat jakým způsobem jsem vybíral ostatní barvy palety (chips, avatar apod.). Nemusím sepisovat každou barvu, ale říct přibližně jaké barvy jsem vybral a jaký tool jsem použil pro zajištění dobré viditelnosti a accessibility (takový to že to má dobrý kontrast). U barev avatarů napsat, že jsem pro pozadí například volil pastelové barvy, na kterých jsou ostatní barvy dobře viditelné.    

% Review
\subsection{Výběr písma}

% Review
\subsection{Výběr zaoblení a rozložení komponentů}

% Review
\subsection{Výběr animací a efektů}

